% math-coding.tex — main file for the formal model.

\documentclass[11pt,a4paper]{article}

% preamble.tex — packages and macros for math-coding v2.0-Y formal model.
% This file is \input{} from math-coding.tex; it must NOT contain
% \documentclass.

\usepackage[utf8]{inputenc}
\usepackage[T1]{fontenc}
\usepackage[english]{babel}

\usepackage{amsmath}
\usepackage{amssymb}
\usepackage{amsthm}

\usepackage{hyperref}
\hypersetup{
  colorlinks=true,
  linkcolor=blue,
  citecolor=blue,
  urlcolor=blue,
  pdftitle={math-coding v2.0-Y: formal model},
  pdfauthor={math-coding contributors}
}

\theoremstyle{plain}
\newtheorem{theorem}{Theorem}
\newtheorem{lemma}[theorem]{Lemma}
\newtheorem{proposition}[theorem]{Proposition}
\newtheorem{corollary}[theorem]{Corollary}

\theoremstyle{definition}
\newtheorem{definition}{Definition}

\newcommand{\Dec}{\mathsf{Decision}}
\newcommand{\Prop}{\mathsf{Proposition}}
\newcommand{\Code}{\mathsf{Code}}
\newcommand{\Witness}{\mathsf{Witness}}
\newcommand{\Actor}{\mathsf{Actor}}
\newcommand{\Reg}{\mathsf{Register}}
\newcommand{\State}{\mathsf{State}}
\newcommand{\Lifecycle}{\mathsf{Lifecycle}}
\newcommand{\Verdict}{\mathsf{Verdict}}
\newcommand{\Conf}{\mathsf{Confidence}}
\newcommand{\Ben}{\mathsf{Beneficiary}}

\begin{document}

\title{math-coding v2.0-Y: Formal Model}
\author{math-coding contributors}
\date{\today}

\maketitle

\begin{abstract}
math-coding is a discipline for recording architectural decisions as
plain-text packets verified by a single static OCaml binary. This
document formalises the type system, the kernel semantics, the five
foundations and three extensions, and the Y-fixed-point property
that makes the convention self-applicable. The OCaml implementation
in \texttt{core/*.ml} mirrors these definitions; any divergence is
recorded as a \texttt{KNOWN DIVERGENCE} comment with an issue
reference.
\end{abstract}

\tableofcontents

% syntax.tex — type definitions for math-coding v2.0-Y.
% Companion to MANIFESTO.md; canonical formal model.

\section{Syntax}

\begin{definition}[Schema version]
A schema version is a SemVer string:
$\mathbb{V} \subseteq \{\, v \in \mathbb{S} \mid v = \mathrm{major}.\mathrm{minor}.\mathrm{patch} \,\}$.
The current canonical version is ``2.0''.
\end{definition}

\begin{definition}[Actor]
An actor identifies who originated a decision:
$$\Actor ::= \mathrm{Human} \mid \mathrm{Agent} \mid \mathrm{System}$$
\end{definition}

\begin{definition}[Register]
The epistemic stance of a decision:
$$\Reg ::= \mathrm{Fact} \mid \mathrm{Hypothesis} \mid \mathrm{Judgment} \mid \mathrm{Unknown}$$
\end{definition}

\begin{definition}[State]
The FSM state of a packet:
$$\State ::= \mathrm{Draft} \mid \mathrm{Applied} \mid \mathrm{Reviewed} \mid \mathrm{Retired} \mid \mathrm{Abandoned}$$
\end{definition}

\begin{definition}[Lifecycle]
The computed lifecycle:
$$\Lifecycle ::= \mathrm{Draft} \mid \mathrm{Applied} \mid \mathrm{Drift} \mid \mathrm{Stale}$$
\end{definition}

\begin{definition}[Verdict]
The verdict of a check:
$$\Verdict ::= \mathrm{Pass} \mid \mathrm{Warn} \mid \mathrm{Fail} \mid \mathrm{Skip}$$
\end{definition}

\begin{definition}[Beneficiary]
Who is served by the decision:
$$\Ben ::= \mathrm{User} \mid \mathrm{Developer} \mid \mathrm{Team} \mid \mathrm{FutureSelf} \mid \mathrm{System} \mid \mathrm{Other}(\mathbb{S})$$
\end{definition}

\begin{definition}[Decision]
A decision is the central data type of math-coding. Its schema:
\begin{align*}
\Dec \;:=\; \langle
& \mathtt{schema\_version} : \mathbb{V}, \\
& \mathtt{name}            : \mathbb{S}, \\
& \mathtt{proposition}     : \mathbb{S}, \\
& \mathtt{code}            : \mathbb{S} \cup \{\bot\}, \\
& \mathtt{witness}         : \mathbb{SHA} \cup \{\bot\}, \\
& \mathtt{register}        : \Reg, \\
& \mathtt{state}           : \State, \\
& \mathtt{actor}           : \Actor, \\
& \mathtt{confidence}      : [0, 1], \\
& \mathtt{superseded\_by}  : \mathbb{S} \cup \{\bot\}, \\
& \mathtt{beneficiary}     : \Ben
\rangle
\end{align*}
Field constraints:
\begin{itemize}
\item $\mathtt{schema\_version} = \mathtt{``2.0''}$ (mandatory, $V_7$)
\item $\mathtt{proposition} \neq \mathtt{``''}$ (mandatory, $V_1$)
\item $\mathtt{register} \in \Reg$ (mandatory, $V_3$)
\item $\mathtt{state} \in \State$ (mandatory, $V_4$)
\item $\mathtt{actor} \in \Actor$ (mandatory, $V_5$)
\item $\mathtt{confidence} \in [0, 1]$ (mandatory)
\item $\mathtt{confidence} \models \mathtt{register}$ (consistency, $V_3$)
\end{itemize}
\end{definition}

\begin{definition}[Register-confidence consistency]
The register constrains the confidence value:
\begin{align*}
\mathtt{register} = \mathrm{Fact}       &\;\Rightarrow\; \mathtt{confidence} \in [0.95,\, 1] \\
\mathtt{register} = \mathrm{Hypothesis} &\;\Rightarrow\; \mathtt{confidence} \in (0.5,\, 0.95) \\
\mathtt{register} = \mathrm{Judgment}   &\;\Rightarrow\; \mathtt{confidence} \in \{0,\, 1\} \\
\mathtt{register} = \mathrm{Unknown}    &\;\Rightarrow\; \mathtt{confidence} = 0
\end{align*}
\end{definition}

\begin{definition}[Substrate]
A substrate specifies how a packet's proof is verified:
\begin{align*}
\mathsf{Substrate} ::= \; & \mathsf{None} \\
  \mid\; & \mathsf{Shell}(\mathbb{S}) \\
  \mid\; & \mathsf{Pbt}(\mathbb{S}) \\
  \mid\; & \mathsf{Tla}(\mathbb{S}) \\
  \mid\; & \mathsf{Coq}(\mathbb{S}) \\
  \mid\; & \mathsf{Alloy}(\mathbb{S}) \\
  \mid\; & \mathsf{Bpmn}(\mathbb{S}) \\
  \mid\; & \mathsf{PbtPrism}(\mathbb{S})
\end{align*}
\end{definition}
% semantics.tex — kernel S semantics for math-coding v2.0-Y.

\section{Semantics}

The kernel $S$ maps a decision (plus repo context) to a verdict.

\begin{definition}[Kernel $S$]
$$S : \Dec \times \mathrm{Repo} \to \Verdict$$
$S$ composes six structural predicates V1--V6 plus V7 schema
validation. The verdict is the worst case across all predicates.
\end{definition}

\subsection{Structural predicates}

\begin{proposition}[V1: structure]
$$\forall d.\ d.\mathtt{proposition} \neq \mathtt{``''}$$
If violated, $V_1(d) = \mathrm{Fail}$.
\end{proposition}

\begin{proposition}[V2: lifecycle]
$$\forall d.\ d.\mathtt{state} = \mathrm{Applied} \iff (d.\mathtt{witness} \neq \bot \;\wedge\; \mathrm{prop\_at\_witness}(d) = d.\mathtt{proposition})$$
Drift ($\mathrm{prop\_at\_witness}(d) \neq d.\mathtt{proposition}$) returns
$\mathrm{Drift}$. No witness returns $\mathrm{Draft}$. Returns
$\mathrm{Warn}$ (kernel emits $\mathrm{Warn}$ verdict).
\end{proposition}

\begin{proposition}[V3: register--confidence consistency]
$$\forall d.\ d.\mathtt{confidence} \models d.\mathtt{register}$$
(Defined in \texttt{syntax.tex}, ``Register-confidence consistency''.)
If violated, $V_3(d) = \mathrm{Fail}$.
\end{proposition}

\begin{proposition}[V4: state FSM]
The FSM permits transitions
$\mathrm{Draft} \to \mathrm{Applied}$, $\mathrm{Applied} \to \mathrm{Reviewed}$,
$\mathrm{Applied} \to \mathrm{Retired}$, $\mathrm{Applied} \to \mathrm{Abandoned}$,
$\mathrm{Applied} \to \mathrm{Draft}$ (via supersession).
The transition $\mathrm{Draft} \to \mathrm{Reviewed}$ is forbidden when
$d.\mathtt{witness} = \bot$. If attempted, $V_4(d) = \mathrm{Fail}$.
\end{proposition}

\begin{proposition}[V5: actor discipline]
Three modes via $.mathrc\!:\!\mathtt{SIGNING\_MODE} \in \{\mathrm{strict}, \mathrm{lenient}, \mathrm{off}\}$:
\begin{itemize}
\item $\mathrm{strict}$: $d.\mathtt{actor} = \mathrm{Human} \wedge d.\mathtt{signed\_by} = \bot \Rightarrow V_5(d) = \mathrm{Fail}$
\item $\mathrm{lenient}$: same condition returns $\mathrm{Warn}$
\item $\mathrm{off}$: signing not checked
\end{itemize}
Universal warnings:
\begin{itemize}
\item $d.\mathtt{actor} = \mathrm{Agent} \wedge d.\mathtt{register} = \mathrm{Fact} \Rightarrow \mathrm{Warn}$ (an agent shall not assert fact without evidence)
\item $d.\mathtt{actor} = \mathrm{Agent} \wedge d.\mathtt{state} = \mathrm{Reviewed} \Rightarrow \mathrm{Warn}$ (state Reviewed requires explicit human sign-off)
\item $d.\mathtt{actor} = \mathrm{Human} \wedge d.\mathtt{motivation} = \bot \Rightarrow \mathrm{Warn}$
\end{itemize}
\end{proposition}

\begin{proposition}[V6: supersession SPO]
Supersession $\preceq$ on decisions is a strict partial order:
\begin{itemize}
\item irreflexive: $\neg(d \preceq d)$
\item asymmetric: $d_1 \preceq d_2 \Rightarrow \neg(d_2 \preceq d_1)$
\item transitive: $d_1 \preceq d_2 \wedge d_2 \preceq d_3 \Rightarrow d_1 \preceq d_3$
\end{itemize}
A cycle in the supersession graph returns $V_6(d) = \mathrm{Fail}$.
\end{proposition}

\begin{proposition}[V7: schema version]
$$d.\mathtt{schema\_version} = \mathtt{``2.0''}$$
If violated (unknown version), $V_7(d) = \mathrm{Fail}$. Future
versions trigger best-effort parsing with a warning.
\end{proposition}

\subsection{Composition}

\begin{definition}[Verdict composition]
The overall verdict is the worst case across V1--V7:
$$\mathrm{Verdict}(S, d) = \min(\{V_i(d) \mid i \in 1..7\})$$
where $\min$ is the total order $\mathrm{Fail} < \mathrm{Warn} < \mathrm{Pass} < \mathrm{Skip}$.
\end{definition}

\begin{definition}[Lifecycle computation]
$$L : \Dec \times \mathrm{Repo} \to \Lifecycle$$
\begin{align*}
L(d, r) = \;& \mathrm{Draft}   &&\text{if } d.\mathtt{witness} = \bot \\
           &\;\mid\; \mathrm{Drift} &&\text{if } d.\mathtt{witness} = \text{sha} \;\wedge\; \mathrm{prop}(r, \text{sha}, d.\mathtt{name}) \neq d.\mathtt{proposition} \\
           &\;\mid\; \mathrm{Stale} &&\text{if } d.\mathtt{witness} = \text{sha} \;\wedge\; \mathrm{prop}(r, \text{sha}, d.\mathtt{name}) = d.\mathtt{proposition} \;\wedge\; \text{files changed since sha} \\
           &\;\mid\; \mathrm{Applied} &&\text{otherwise}
\end{align*}
\end{definition}
% foundations.tex — five foundational propositions of math-coding.

\section{Foundations}

Five foundational propositions ground the convention. Each is
instantiated as a packet under \texttt{math/foundations/}.

\subsection{curry-howard: decision is a triple}

\begin{theorem}[curry-howard.triple]
A decision in math-coding is structurally a triple $(\Prop, \Code, \Witness)$,
and the structural correspondence between proposition and code is
verified through witness.
\end{theorem}

\begin{proof}
Curry-Howard isomorphism states that proposition is a type, code is a
term, and inhabitation is a proof. In software engineering, proposition
becomes a statement about a decision, code its realisation, and
witness the git commit binding the proposition to that realisation.
Without the triple, proposition is documentation without consequences:
its change cannot be tracked. With the triple, proposition and code
become a verifiable pair; any code change either updates the
proposition or triggers an explicit supersession.
\end{proof}

\begin{theorem}[Triple in OCaml]
The OCaml record \texttt{decision} mirrors the triple directly:
\begin{verbatim}
type decision = {
  proposition : string;     (* type *)
  code        : string option; (* term *)
  witness     : Sha.t option;  (* inhabitation proof *)
}
\end{verbatim}
\end{theorem}

\subsection{temporal: lifecycle is computed}

\begin{theorem}[temporal.lifecycle]
The lifecycle of a decision $L$ is a function of the decision and git
history; $L$ is computed, not stored.
\end{theorem}

\begin{proof}
Decisions evolve. A proposition true at commit $A$ may be false at commit
$B$. Storing lifecycle as a stored field admits lying: a developer
changes code but forgets to update the \texttt{state} field. Computing
lifecycle from git history makes drift detection automatic and
falsifiable---a lie about state becomes mechanically impossible,
because state is a consequence of history, not an editable field.
\end{proof}

\begin{theorem}[Lifecycle function]
\begin{align*}
L : \Dec \times \mathrm{Repo} &\to \Lifecycle \\
L(d, r) = \;& \textsc{Draft}   &&\text{if } d.\mathtt{witness} = \bot \\
           &\;\mid\; \textsc{Drift} &&\text{if } d.\mathtt{witness} = \text{sha} \;\wedge\; \mathrm{prop}(r, \text{sha}, d.\mathtt{name}) \neq d.\mathtt{proposition} \\
           &\;\mid\; \textsc{Stale} &&\text{if } d.\mathtt{witness} = \text{sha} \;\wedge\; \mathrm{prop}(r, \text{sha}, d.\mathtt{name}) = d.\mathtt{proposition} \;\wedge\; \text{files changed since sha} \\
           &\;\mid\; \textsc{Applied} &&\text{otherwise}
\end{align*}
\end{theorem}

\subsection{constructive: proof is reproducible}

\begin{theorem}[constructive.proof]
A packet with \texttt{substrate: shell} or \texttt{substrate: pbt} requires
reproducible evidence: re-running the command yields the recorded
exit code.
\end{theorem}

\begin{proof}
Classical logic permits ``proven'' without construction. Constructive
logic requires a term---a program that normalises to the sought value.
In software engineering, the term is a re-runnable command: test,
benchmark, type-check. If the result is not reproducible, the claim is
not proven---it is a hypothesis. This is the only epistemic marker
that triggers runtime action; \texttt{register: fact}, \texttt{hypothesis},
\texttt{judgment}, \texttt{unknown} are declarative stances.
\end{proof}

\begin{theorem}[Reproducible definition]
$$\mathrm{Reproducible} \triangleq \mathrm{Re\!-\!run}(\mathrm{command}) = \mathrm{recorded\_exit}$$
\end{theorem}

\subsection{categorical: supersession is SPO}

\begin{theorem}[categorical.supersession-spo]
Supersession $\preceq$ on decisions is a strict partial order:
irreflexive, asymmetric, transitive.
\end{theorem}

\begin{proof}
Decisions evolve. When a proposition changes, the old decision is not
edited---it is superseded by a new one. The relation ``supersedes''
must be a strict partial order to preserve meaning.
Irreflexivity prevents trivial self-supersession. Asymmetry excludes
cycles. Transitivity composes chains into lineages. Without these
properties, the history of decisions becomes inconsistent: a
reviewer cannot tell which decision is current.
\end{proof}

\begin{theorem}[SPO axioms]
\begin{align*}
\forall a.\ &\neg(a \preceq a)                              &&\text{(irreflexive)} \\
\forall a, b.\ &(a \preceq b) \Rightarrow \neg(b \preceq a) &&\text{(asymmetric)} \\
\forall a, b, c.\ &(a \preceq b) \wedge (b \preceq c) \Rightarrow (a \preceq c) &&\text{(transitive)}
\end{align*}
\end{theorem}

\subsection{motivation: register and why}

\begin{theorem}[motivation.register-and-why]
Every packet declares an epistemic register
(\textsc{Fact}/\textsc{Hypothesis}/\textsc{Judgment}/\textsc{Unknown}) and
a numerical confidence. A convention without motivation is ceremony.
\end{theorem}

\begin{proof}
Decisions are taken with different degrees of confidence. A fact
requires reproduction; a hypothesis is open to refutation; a judgment
is an explicit indication of potential harm (reversibility,
mitigation); an unknown is an honest admission of a boundary. Without
explicit register, epistemic honesty degrades into self-deception:
a developer fills ``fact'' without evidence; an agent fills
``fact'' without verification. Without motivation, proposition
becomes a slogan beyond critique. This is axiom A1 (Care) of
math-coding v0.854, restored in v2.0-Y.
\end{proof}

\begin{theorem}[Register-confidence consistency]
The register constrains the confidence value:
\begin{align*}
\mathtt{register} = \textsc{Fact}       &\;\Rightarrow\; \mathtt{confidence} \in [0.95,\, 1] \\
\mathtt{register} = \textsc{Hypothesis} &\;\Rightarrow\; \mathtt{confidence} \in (0.5,\, 0.95) \\
\mathtt{register} = \textsc{Judgment}   &\;\Rightarrow\; \mathtt{confidence} \in \{0,\, 1\} \\
\mathtt{register} = \textsc{Unknown}    &\;\Rightarrow\; \mathtt{confidence} = 0
\end{align*}
The \texttt{judgment} register additionally obligates a non-empty
\texttt{\#\# Why} section in the packet body, declaring reversibility
and mitigation.
\end{theorem}
% extensions.tex — three extension propositions of math-coding.

\section{Extensions}

Three extensions add domain-specific obligations to the foundation
schema. Each is instantiated as a packet under
\texttt{math/extensions/}.

\subsection{process-fsm: state machine with one forbidden transition}

\begin{theorem}[process-fsm.fsm]
A decision exists in one of five states
$\textsc{Draft} \mid \textsc{Applied} \mid \textsc{Reviewed} \mid \textsc{Retired} \mid \textsc{Abandoned}$.
The single forbidden transition is $\textsc{Draft} \to \textsc{Reviewed}$
without $\Witness$.
\end{theorem}

\begin{proof}
Without an explicit state machine, a packet remains ``applied'' by
default and review loses meaning. With a five-state machine, each
decision traverses an explicit path: from proposition (Draft) to
inhabitation (Applied) to human review (Reviewed), with possible
exit through Retired or Abandoned. The forbidden transition
$\textsc{Draft} \to \textsc{Reviewed}$ without $\Witness$ closes a
trivial bypass: one cannot declare a decision reviewed without
passing through inhabitation via git fix. This is the minimal
sufficient restriction: one rule closes all bypasses.
\end{proof}

\begin{theorem}[process-fsm.forbidden-transition]
$$\forall d.\ d.\mathtt{state} = \textsc{Reviewed} \;\wedge\; d.\mathtt{witness} = \bot \;\Rightarrow\; S(d) = \textsc{Fail}$$
The kernel $\mathrm{V_4}$ detects this and refuses the packet.
\end{theorem}

\subsection{dialectic-tas: thesis, antithesis, synthesis required}

\begin{theorem}[dialectic-tas.required-sections]
A packet with $\Reg = \textsc{Judgment}$ or $\Actor = \textsc{Human}$
requires the body sections \texttt{\#\# Thesis}, \texttt{\#\# Antithesis},
\texttt{\#\# Synthesis}. The structural form of dialectics is mandatory
for human judgments.
\end{theorem}

\begin{proof}
Human judgment without explicit consideration of counter-arguments
slips into self-deception. The dialectic tradition shows that
synthesis is reachable only through thesis and antithesis; collapsing
this structure turns judgment into declaration. In software
engineering, judgment about trust, security, or reversibility must
contain an explicit antithesis---the strongest counter-argument
against the decision taken. Without it, the proposition is a slogan
and the reviewer cannot distinguish considered decision from
impulsive one.
\end{proof}

\subsection{actor-discipline: signed commits fix author}

\begin{theorem}[actor-discipline.signed-commits]
A decision records its author through commit signatures. The signing
mode ($\textsc{strict} / \textsc{lenient} / \textsc{off}$) is determined
by the project through \texttt{.mathrc} and checked by the kernel.
\end{theorem}

\begin{proof}
A decision without author carries no responsibility. In distributed
development (human, agent, agentic system), commit signature is the
only cryptographically verifiable way to bind a proposition to a
subject. Without an explicit signing mode, the project ends up in a
situation where ``everyone signs, no one checks''---or vice versa,
``no one signs, authorship is untraceable''. Three modes---strict,
lenient, off---cover three classes of projects: critical
infrastructure, ordinary development, prototyping.
\end{proof}

\begin{theorem}[Signing modes]
\begin{itemize}
\item \textsc{strict}: every commit of the packet has a valid GPG/SSH
      signature. Kernel rejects the packet without signature
      ($\textsc{Fail}$).
\item \textsc{lenient}: the \texttt{amend} commit has a signature;
      others need not. Kernel returns $\textsc{Warn}$ if \texttt{amend}
      lacks signature.
\item \textsc{off}: signatures not required; kernel does not check.
\end{itemize}
\end{theorem}
% theorems.tex — derived theorems and the Y-fixed-point.

\section{Theorems}

\subsection{Y-fixed-point: math-coding applies to itself}

\begin{theorem}[y-fixed-point.math-coding]
math-coding is a Y-fixed point:
$$Y(\mathrm{math\text{-}coding}) = \mathrm{math\text{-}coding}$$
The convention, applied to itself, yields itself.
\end{theorem}

\begin{proof}
math-coding is a function $\mathcal{M} : \mathrm{Project} \to
\mathcal{P}(\Dec)$. Applied to itself, it produces:
\begin{enumerate}
\item the five foundations as ordinary packets
      (\texttt{curry-howard}, \texttt{temporal}, \texttt{constructive},
      \texttt{categorical}, \texttt{motivation}),
\item the three extensions as ordinary packets
      (\texttt{process-fsm}, \texttt{dialectic-tas},
      \texttt{actor-discipline}),
\item the kernel $S$ that type-checks those packets using the same
      code that checks user packets.
\end{enumerate}
All three are themselves packets verified by the same kernel. There
is no special path for the foundations; this is the structural
expression of the Y-fixed-point property. If any proposition in a
foundation packet is false, the kernel detects drift in that
foundation just as it does for a user packet.
\end{proof}

\subsection{Drift detection}

\begin{theorem}[drift.detection]
$$L(d, r) = \mathrm{Drift} \iff d.\mathtt{witness} = \text{sha} \;\wedge\; \mathrm{prop}(r, \text{sha}, d.\mathtt{name}) \neq d.\mathtt{proposition}$$
That is, drift occurs exactly when the proposition at the witness
commit differs from the current proposition.
\end{theorem}

\begin{proof}
If $\mathrm{prop}(r, \text{sha}, d.\mathtt{name}) = d.\mathtt{proposition}$,
the proposition is consistent with history and $L = \mathrm{Applied}$
(possibly $\mathrm{Stale}$ if files changed). If they differ, the
proposition has been modified without a supersession, which is the
definition of drift. The kernel implements this comparison through
\texttt{git show sha:math/name/packet.md}.
\end{proof}

\subsection{SPO transitivity preserves lineage}

\begin{theorem}[spo-transitivity]
If $d_1 \preceq d_2$ and $d_2 \preceq d_3$, then $d_1 \preceq d_3$
through the composition of two supersession chains.
\end{theorem}

\begin{proof}
By definition of transitivity in
\texttt{categorical.supersession-spo}, the composition of two
supersession links is a supersession link. The kernel preserves
chains through the \texttt{mathc graph} command, which renders the
Hasse diagram of the partial order.
\end{proof}

\subsection{Kernel uniformity}

\begin{theorem}[kernel-uniformity]
The kernel $S$ applies uniformly to all decisions, including the
foundations and extensions. No special path exists for axiomatic
packets.
\end{theorem}

\begin{proof}
\texttt{core/check.ml} iterates over \texttt{math/} recursively and
applies \texttt{Check.check : decision -> verdict} to every packet
without distinction. The foundations satisfy the schema (V1--V7)
the same way user packets do. This is the operational expression of
the Y-fixed-point: the kernel that verifies packets also verifies
the description of the kernel.
\end{proof}

\section*{Colophon}

This document was typeset with \texttt{pdflatex} twice for cross-
references. The source lives in \texttt{math/modeling/} alongside
the OCaml source in \texttt{core/}. Both layers are kept in sync
through manual review and the \texttt{KNOWN DIVERGENCE} convention.

\end{document}